\section{Further research questions}

\begin{enumerate}
\item Can the fixed-probe opening rule be represented in the smallest
arity-based certificate language without losing the source savings to
normalization and scope bookkeeping?  The present 508-byte program belongs
to the inherited calculus, not to this new rule basis.
\item Can the entire substitution theorem for deterministic opening be
formalized directly in Lean?  The interesting case is a substitution that
creates an extractable equality where the old opening used a fresh
parameter.  Its branch must be instantiated at the newly determined term.
\item Can a suitably designed canonical syntax expose both the zero row
and first row without full normalization, while retaining all current
formula round trips?  A syntactic shortcut must justify its whole-domain
interpretation and count its decoder.
\item Is there a complete presentation improving the token, AST, and byte
coordinates simultaneously under one fixed input contract?  The present
coordinate leaders remain different programs.
\item Can proof translations between the new presentation and the raw
native core be given as fully checked certificate compilers, including
normalization traces and the introduction-mode repair?  The current
universal result is a written metatheorem, not a claimed general executable
reverse compiler.
\item What are useful lower bounds under a genuinely fixed implementation
language and public theorem interface?  The observed corpus minima alone
provide no such lower bounds.
\end{enumerate}

